% Files of the document named like tex4ht's .4ht files, read in the
% preamble: glossary, abstract, appendix, and index with \input, and
% background with \InputIfFileExists. Their names end in .ltx, so that
% they are not fixtures. A glossary.tex in a subfolder, read with the
% folder in its name after glossary.ltx, is another file. As in the PDF,
% each file's text is in the page, and the abstract and the appendix keep
% article's form. A file read with the primitive \input, one found through
% \input@path, and one read after \begin{document} are read the same way.
% makeidx's index keeps its entry. The glossary and abstract files name
% themselves with \ProvidesFile under another name than their file's.
\documentclass{article}
\usepackage{makeidx}
\makeindex
\input{glossary.ltx}
% A file of test-inputname.tex, named like glossary.4ht, read with its
% folder in its name after the glossary.ltx beside the document.
\newcommand\ownsubglossary{Text from the subfolder glossary file.}

\input{abstract.ltx}
\input{appendix.ltx}
\input{index.ltx}
\InputIfFileExists{background.ltx}{}{\newcommand\ownbackground{No background file.}}
\input reading.ltx
\makeatletter
\def\input@path{{test-inputname/}}
\makeatother
\input{proof}
\begin{document}
\begin{abstract}
The abstract of the document.
\end{abstract}
\section{Files}
A ring\index{ring}.

Glossary: \ownglossary

Subfolder: \ownsubglossary

Abstract: \ownabstract

Appendix: \ownappendix

Index: \ownindex

Background: \ownbackground

Reading: \ownreading

Proof: \ownproof

Errata: \input{errata.ltx}
\appendix
\section{Tables}
The appendix of the document.
\printindex
\end{document}
